% Generated draft: requires mathematical and source review.
\documentclass[11pt]{article}
\usepackage[margin=25mm]{geometry}
\usepackage{fontspec}
\setmainfont{Latin Modern Roman}
\usepackage{amsmath,amssymb,mathtools,unicode-math}
\setmathfont{Latin Modern Math}
\usepackage{enumitem,needspace,xurl,hyperref}
\setlength{\parindent}{0pt}
\setlength{\parskip}{5pt}
\setlength{\emergencystretch}{3em}
\newcommand{\N}{\mathbb{N}}
\newcommand{\Z}{\mathbb{Z}}
\newcommand{\Q}{\mathbb{Q}}
\newcommand{\R}{\mathbb{R}}
\newcommand{\C}{\mathbb{C}}
\newcommand{\F}{\mathbb{F}}
\newcommand{\A}{\mathbb{A}}
\newcommand{\PP}{\mathbb{P}}
\newcommand{\E}{\mathbb{E}}
\newcommand{\eps}{\varepsilon}
\newcommand{\dd}{\,\mathrm{d}}
\begin{document}
\section*{UnsolvedMath ID: 30000516}
\textbf{Source problem code:} catalogue entry
\subsection{Definitions and mathematical statement}
For a Beauville surface $S$, determine a field of definition of $S$ and the $\operatorname{Gal}(\overline{\Q}/\Q)$-orbit of its moduli point $[S]$.

The source gives no presentation of $S$. In the usual terminology a Beauville surface is a rigid surface isogenous to a product of curves of genus at least two; in an unmixed presentation one writes $S=(C_1\times C_2)/G$ with a free diagonal finite-group action and the relevant rigidity data. The extraction does not specify an unmixed or mixed case, curves, group, or ramification data, so the unmixed presentation is only explanatory and is not imposed as a source hypothesis.

Fixing the algebraic-closure and complex-field identifications needed to interpret the question, a field of definition $K\subseteq\overline{\Q}$ is one for which there exists a surface $S_K$ with
\[
 S_K\times_K\overline{\Q}\cong S.
\]
If the given surface is initially viewed over $\C$, the right-hand side is interpreted after the chosen embedding into $\C$ and requires the algebraic model intended by the source. The question asks for a field of definition, not explicitly a smallest such field or a descent theorem to the field of moduli.

For $\sigma\in\operatorname{Gal}(\overline{\Q}/\Q)$, let $S^\sigma$ be the conjugate surface. The orbit to determine is
\[
 \{[S^\sigma]:\sigma\in\operatorname{Gal}(\overline{\Q}/\Q)\}.
\]
This concerns isomorphism classes of surfaces as moduli points, unless markings or covering data are expressly included. The stabilizer is $\{\sigma:S^\sigma\cong S\}$; its fixed field describes the field-of-moduli notion, which is not automatically identified here with a field of definition.

A concrete answer requires the particular Beauville surface or a precisely specified input family and the intended moduli equivalence. The dated assessment says the definition of $S$ is omitted. Naming the general surface class does not recover those missing instance data, and no surface, group, field, or Galois orbit is invented.
\subsection{Short English statement}
The question asks for an algebraic field model and all Galois conjugate moduli classes of a Beauville surface. The specific surface is missing, and the field of moduli must be kept distinct from an actual field of definition.
\subsection{Sources}
\begin{itemize}
\item Local UnsolvedMath registry, permanent ID 30000516, code catalogue entry, ``Fields of Definition and Galois Orbits of Beauville Surfaces.'' Proposed year 2006; proposer unspecified.
\item \emph{Classical Algebraic Geometry} (2007), Oberwolfach Reports, pp.\ 1615--1662, DOI 10.4171/owr/2006/27; \url{https://doi.org/10.4171/owr/2006/27}. The displayed report year 2007 differs from the DOI and proposed year 2006.
\item Statement status corrected\_verified; the supplied verification says official-report content was checked and notation, context, and question form normalized without changing the source claim.
\item Literature triage checked 2026-08-21: open. It says the missing definition of $S$ prevents identifying a field of definition or Galois orbit, so no determinate resolution is available. Created 2026-08-25; updated 2026-09-16. This is a dated formulation assessment requiring source verification, not a status theorem about every Beauville surface.
\item Registry curation by Ulam AI / UnsolvedMath Contributors, CC BY 4.0; the underlying report retains its supplied provenance. Unreviewed statement-only draft.
\end{itemize}
\end{document}
